\section{Introducción}

El presente informe documenta el desarrollo del Proyecto Grupal I del curso de Big Data, cuyo objetivo es aplicar de forma práctica los conceptos de almacenamiento distribuido, procesamiento paralelo y programación en la nube sobre un conjunto de datos real.

Para ello, el grupo trabajó con un dataset de transacciones de comercio electrónico (\emph{Online Retail II}) que supera el millón de registros, y construyó una arquitectura completa sobre Google Cloud Platform (GCP): un data lake en Cloud Storage, un clúster de Google Cloud Dataproc, y un pipeline de procesamiento implementado en cuatro frameworks distintos (Polars, Dask, Modin y Apache Spark), además de tres programas de Hadoop MapReduce ejecutados sobre HDFS.

El objetivo no fue únicamente demostrar que cada herramienta ``corre'', sino resolver, con las cuatro, el mismo conjunto de doce consultas de negocio sobre los datos —limpieza, deduplicación, transformación, filtrado, agregaciones, agrupaciones, ordenamiento y cálculo de métricas— para poder comparar de forma directa cómo cada framework resuelve exactamente el mismo problema. Toda la infraestructura de GCP (bucket y clúster) se define como código mediante Terraform, lo que permite recrearla o destruirla de forma reproducible.

Las secciones siguientes describen, en orden: el dataset elegido y su calidad de datos, la implementación del almacenamiento y la arquitectura sobre GCP, el procesamiento distribuido con los cuatro frameworks, las operaciones concretas realizadas sobre los datos, y finalmente la ejecución de Hadoop MapReduce sobre el clúster.
